\documentclass[11pt,a4paper]{article}
\usepackage[main=italian,provide=*]{babel}
\usepackage[utf8]{inputenc}
\usepackage[T1]{fontenc}
\usepackage{amsmath,amssymb}
\usepackage{geometry}
\geometry{margin=2.7cm}
\usepackage{booktabs}
\usepackage{array}
\usepackage{caption}
\usepackage{seqsplit}
\usepackage[colorlinks=true,linkcolor=blue,citecolor=blue,urlcolor=blue]{hyperref}

\newcommand{\ket}[1]{\lvert #1 \rangle}
\newcommand{\bra}[1]{\langle #1 \rvert}

\title{Quantum simulation del trimero a catena aperta: applicazione della teoria\\
all'implementazione Qiskit\\[6pt]
\large Dalla teoria a tre livelli al modulo \texttt{\seqsplit{trotter\_trimero\_catena.py}} --- procedura commentata}
\author{Samuele \\ Relatore: Prof. Alessandro Chiesa}
\date{}

\begin{document}
\maketitle

\begin{abstract}
\noindent Questo documento ricostruisce, passo per passo, come la teoria
raccolta in \texttt{\seqsplit{quantum\_simulation\_trimero\_catena\_teoria.tex}} sia
stata applicata nella costruzione del modulo
\texttt{\seqsplit{trotter\_trimero\_catena.py}} e nell'analisi che ne è seguita:
ricerca del punto di lavoro, convergenza in $N$, il risultato originale
dell'alternanza dell'ordine dei bond, studio di compattezza del circuito.
Segue lo stesso schema di
\texttt{\seqsplit{quantum\_simulation\_trimero\_anello\_applicazione.tex}}, con una
differenza di partenza da tenere presente: qui il punto di lavoro non ha
(ancora) una derivazione fisica principiata, resta uno scan numerico
verificato robusto -- stesso status provvisorio di $R_0$ (trimero anello).
\end{abstract}

\tableofcontents

\section{Mappa dei riferimenti}

\begin{center}
\begin{tabular}{@{}p{4.6cm}p{9.6cm}@{}}
\toprule
Documento & Stato e uso \\
\midrule
\texttt{\seqsplit{quantum\_simulation\_trimero\_catena\_teoria.tex}} & Derivazione completa dei blocchi esatti e dell'errore a tre livelli. Riferimento primario di questo documento. \\
\texttt{\seqsplit{quantum\_simulation\_trimero\_anello\_teoria.tex}} / \texttt{\_applicazione.tex} & Template metodologico diretto: stessa struttura di derivazione, qui con due legami invece di tre. \\
\texttt{\seqsplit{log\_decisioni.md}} & Cronologia: correzione del non-sequitur ereditato dall'anello, scoperta dell'alternanza, punto $S_0$ adottato. \\
\texttt{\seqsplit{trimer\_chain\_exact.py}} & Sorgente unica dell'Hamiltoniana, riusata direttamente da \texttt{\seqsplit{trotter\_trimero\_catena.py}}. \\
\texttt{\seqsplit{analisi\_dm\_trimero\_catena.tex}} & Derivazione del vincolo $D_{12}=-D_{23}$ da $P_{13}$, e della cautela su $b_c$ fisso vs.\ minimo vero del gap. \\
\texttt{\seqsplit{trimero\_catena\_circuito\_compatto.tex}} & Studio di compattezza (\S8 di questo documento ne riporta solo il riepilogo). \\
\bottomrule
\end{tabular}
\end{center}

\section{Punto di partenza e decisione metodologica}

\subsection{Perché $H_{ex}$ non è un blocco singolo, anche con due soli legami}

$H_{ex}=J\vec\sigma_1\!\cdot\!\vec\sigma_2+J\vec\sigma_2\!\cdot\!\vec\sigma_3$ condivide il qubit 2 fra i suoi due termini. Stessa scelta metodologica dell'anello: \textbf{strategia (a)}, Trotter interno con gate fisicamente etichettati, non sintesi numerica dell'operatore $8\times8$ (Sezione~7 di questo documento ne riporta il riepilogo per la catena). Introduce il livello 1, presente identico nello spirito a quello dell'anello ma con un solo commutatore da gestire, non tre.

\subsection{Convenzione: differenze esplicite dall'anello}

\begin{center}
\begin{tabular}{@{}p{3.6cm}p{4.4cm}p{4.4cm}@{}}
\toprule
& \texttt{\seqsplit{trotter\_trimero\_anello.py}} & \texttt{\seqsplit{trotter\_trimero\_catena.py}} (qui) \\
\midrule
Legami di scambio & 3 ($J$ base, $J'$ lati) & 2 ($J$ unico) \\
Gradi di libertà DM & 2 opzioni (A/B) & 1 (fissato da $P_{13}$) \\
Normalizzazione DM & $D_{ij}=rJ_{ij}$ & $D_{12}{=}{+}D,\,D_{23}{=}{-}D$ (assoluto) \\
Ordine dei bond & fisso (nessuna alternanza pulita) & alternabile fra passi \\
\bottomrule
\end{tabular}
\end{center}

\textbf{Attenzione}: la riga sulla normalizzazione del DM non è un dettaglio cosmetico. Usare $D_{ij}=D\cdot J_{ij}$ (convenzione anello) al posto di $D_{12}{=}{+}D,D_{23}{=}{-}D$ (convenzione catena) farebbe divergere circuito e benchmark esatto di un fattore $J$, silenziosamente -- verificato che il modulo segue \texttt{\seqsplit{trimer\_chain\_exact.py}::dm\_term}, non l'anello.

\section{Correzione di un non-sequitur ereditato, non propagato}

La documentazione del modulo Trotter dell'anello giustifica il livello 0 esterno affermando che $H_{ex}$ non contribuisce all'errore «perché commuta col campo» -- un non-sequitur: $[S_z^{tot},H_{ex}]=0$ non implica $[H_{ex},H_{DM}]=0$. Verificato numericamente ($J{=}1,b{=}0.7,D{=}0.3$): $\|[H_{ex},H_{DM}]\|=10.182338$ contro $\|[S_z^{tot},H_{DM}]\|=3.394113$ -- $H_{ex}$ domina, non è trascurabile. \textbf{Decisione}: il filone anello resta chiuso, nessuna modifica a \texttt{\seqsplit{trotter\_trimero\_anello.py}} (l'implementazione lì è comunque corretta, solo la frase esplicativa non lo è); la teoria della catena (\texttt{\seqsplit{quantum\_simulation\_trimero\_catena\_teoria.tex}} \S8) è scritta da zero con la giustificazione corretta, senza ereditare l'errore.

\section{La ricerca del punto di lavoro $S_0$}

\subsection{Criterio, identico nello spirito all'anello}

Stesso criterio a due indicatori (escursione picco-picco, rapporto $a_2/a_1$ fra i due modi spettrali dominanti nella decomposizione di Bohr di $\langle S_z^{tot}\rangle(t)$), applicato però in una regione fisicamente motivata anziché su una griglia cieca: attorno al campo critico $b_c=3J$, dove per $D=0$ si ha un crossing esatto (teorema di Kramers, $N$ dispari) e il DM apre un vero anti-crossing (\texttt{\seqsplit{analisi\_dm\_trimero\_catena.tex}}).

\subsection{Una trappola nota, verificata anche qui}

A $b=b_c$ fisso il rapporto $a_2/a_1\to1$ per ogni $D$ testato -- ma il punto di lavoro fisicamente corretto è il \textbf{minimo vero del gap} (funzione \texttt{dm\_min\_gap} di \texttt{\seqsplit{trimer\_chain\_exact.py}}), non $b_c$ fisso: stesso errore metodologico da evitare già documentato per l'anello, verificato qui esplicitamente anziché assunto per analogia.

\subsection{Punto adottato}

\begin{equation*}
\boxed{S_0 = \big(J{=}1,\ D{=}0.3,\ b{=}3.0073414\big)}, \qquad \text{gap}=1.468777.
\end{equation*}
Due modi spettrali dominanti quasi degeneri in ampiezza ($a\approx2.0$, gap $0.847$ e $1.469$): $\langle S_z^{tot}\rangle(t)$ mostra un battimento di periodo $\approx10.1$, non un singolo coseno -- escursione picco-picco $1.34$ su $\ket{010}$. \textbf{Dichiarato esplicitamente provvisorio}, stesso status di $R_0$ (trimero anello): scan verificato, non derivazione fisica principiata.

\section{Una scoperta metodologica di servizio: $\ket{000}$ è cieco al livello 1}

$\ket{000}$ è autostato di \emph{ogni} bond separatamente ($h_{ij}\ket{00}=+\ket{00}$, verificato $\|h_{12}\ket{000}-\ket{000}\|=0$): tutti i fattori Trotter di $H_{ex}$ agiscono su di esso come fasi globali e commutano. Un test di convergenza su $\ket{000}$ a $D=0$ dà infedeltà $\sim10^{-14}$ per \emph{qualunque} $N$ -- un test vuoto. \texttt{\seqsplit{trotter\_trimero\_anello.py}} usa $\ket{000}$ nel suo self-test di convergenza (misura quindi il solo livello 2, non il livello 1) -- non corretto (filone anello chiuso), ma non replicato: qui si usa $\ket{010}$ per ogni analisi di convergenza.

\section{Convergenza in $N$ e alternanza dell'ordine: risultati}

\subsection{Convergenza a $S_0$, ordine fisso}

Su $\ket{010}$, $t=20$: infedeltà $6.32\times10^{-6}$ a $N=8000$, scaling pulito $O(1/N^2)$ (rapporto $\times4.0$ costante raddoppiando $N$), coerente con la teoria (\texttt{\seqsplit{quantum\_simulation\_trimero\_catena\_teoria.tex}} \S9.5: infedeltà $\propto\|\Delta U\|^2\propto(1/N)^2$).

\subsection{Il risultato originale: alternanza dell'ordine dei bond}

Poiché il segno del termine $\chi$ di livello 1 dipende dall'ordine dei due bond (teoria \S4.4), alternarlo fra passi consecutivi cancella quel termine al leading order, a parità esatta di gate -- una possibilità che l'anello (tre bond, non due) non ha in forma altrettanto pulita.

\begin{center}
\small
\begin{tabular}{@{}p{3.4cm}p{2.6cm}p{3.6cm}p{3.6cm}@{}}
\toprule
regime & fisso & alternato & guadagno \\
\midrule
$D=0$, stato casuale & $O(1/N^2)$, $\times4.0$ & $O(1/N^4)$, $\times16.0$ & $34.5\times$ a $N{=}80$, cresce come $N^2$ \\
$D=0.3$, stato casuale & $O(1/N^2)$, $\times4.0$ & $\times9.9\to\times5.1$ (degrada) & $5.7\times$ a $N{=}80$, satura \\
$D=0.3$, punto $S_0$, $\ket{010}$ & $O(1/N^2)$, $\times4.0$ & rapporto $\times5.08$ (non $\times16$) & $25.2\times$ a $N{=}8000$ \\
\bottomrule
\end{tabular}
\end{center}

\textbf{Lettura onesta}: il guadagno di un ordine intero ($O(1/N^4)$) vale \emph{solo} per il livello 1 isolato ($D=0$). Con il DM acceso, l'alternanza cancella il termine chirale ma lascia intatti il livello 2 e i termini incrociati del livello 0 esterno, che restano $O(1/N^2)$ e tornano a dominare per $N$ grande: il guadagno è reale ma parziale e satura, riconfermato sia su stato casuale sia sul punto fisico $S_0$.

\section{Studio di compattezza del circuito}

Ripetuto per la catena l'esperimento già fatto per anello e dimero. A differenza dell'anello, qui non è stato necessario invocare l'assenza di un teorema di compressione chiuso per 3 qubit come limite \emph{concettuale}: il punto resta comunque solo empirico (nessuna garanzia di ottimalità del transpiler), riportato per esteso in \texttt{\seqsplit{trimero\_catena\_circuito\_compatto.tex}}; qui solo il riepilogo:

\begin{itemize}
\item un singolo passo esplicito (10 gate a due qubit nativi a $S_0$) si comprime a \textbf{19 CNOT}, fedeltà $1.0$;
\item il prodotto a $N$ passi si comprime a \textbf{19 CNOT indipendentemente da $N$} (verificato fino a $N=100$, dove il circuito esplicito avrebbe $1000$ gate a due qubit) -- stesso fenomeno di anello e dimero;
\item conclusione invariata: non usare il compresso per lo studio Trotter-vs-rumore della Parte 2.
\end{itemize}

\section{Aperture per il seguito}

\begin{enumerate}
\item \textbf{Derivazione principiata di $S_0$}: sostituire lo scan con una derivazione dalla struttura degli 8 livelli della catena (multipletti $A',B',C'$ di Kambe), analogamente a quanto pianificato per l'anello.
\item \textbf{Parte 2 (rumore)}: il costo in gate per passo (10 nativi a $S_0$) e la scala di $N$ trovata sono il punto di aggancio diretto; l'alternanza dell'ordine dei bond è un candidato concreto per ridurre $N$ a parità di rumore nel regime a DM debole.
\item \textbf{Correlazioni dinamiche su $N=3$, catena}: da estendere con l'ancilla (Hadamard test), riusando \texttt{W-2qC.K2} (Fase 5) come stato di preparazione.
\end{enumerate}

\section*{Riferimenti bibliografici}
\addcontentsline{toc}{section}{Riferimenti bibliografici}

\begin{itemize}
\item H.\,F.\ Trotter, \emph{On the product of semi-groups of operators}, Proc.\ Am.\ Math.\ Soc.\ \textbf{10}, 545--551 (1959).
\item M.\ Suzuki, \emph{Generalized Trotter's formula...}, Commun.\ Math.\ Phys.\ \textbf{51}, 183--190 (1976).
\item S.\ Lloyd, \emph{Universal quantum simulators}, Science \textbf{273}, 1073--1078 (1996).
\item M.\ Crippa et al., \emph{Simulating Static and Dynamic Properties of Magnetic Molecules with Prototype Quantum Computers}, Magnetochemistry \textbf{7}, 117 (2021).
\end{itemize}

\appendix
\section{File prodotti in questa fase del lavoro}

\begin{itemize}
\item \texttt{\seqsplit{trotter\_trimero\_catena.py}} -- modulo Trotter, sette self-test superati a precisione macchina.
\item \texttt{\seqsplit{quantum\_simulation\_trimero\_catena\_teoria.tex}} -- derivazione teorica completa a tre livelli.
\item \texttt{\seqsplit{trimero\_catena\_circuito\_compatto.tex}} -- studio di compattezza del circuito.
\item \texttt{\seqsplit{quantum\_simulation\_trimero\_catena\_trotter.ipynb}} -- notebook di lavoro, eseguito per intero.
\item \texttt{\seqsplit{log\_decisioni.md}} -- aggiornato con la correzione del non-sequitur, il punto $S_0$, l'alternanza.
\end{itemize}

\end{document}
